% The two convexity inequalities, drawn on the same parabola.
%
%   figures/tikz/ra-convexity.tex
%       -> media/figures/ra-convexity.{png,svg}
%
% Lecture 13 (Real Analysis Review), Convexity. Left: inequality (13-1), the
% chord lies above the graph (Alex's own scan picture, relabelled as in the
% previous inline figure). Right: inequality (13-2), the tangent line lies
% below the graph.
%
% GEOMETRY (x unit 1.0 cm, y unit 0.9 cm). phi(x) = 0.8 x^2, domain -1.7..1.7.
%  Left: x_a = -1.5, x_b = 1.0, theta = 1/2, so x_c = -0.25.
%   phi(x_a) = 1.8, phi(x_b) = 0.8, chord value at x_c = 1.3, phi(x_c) = 0.05.
%  Right: expansion point x = -0.5: phi(x) = 0.2, phi'(x) = 1.6 x = -0.8, so the
%   tangent is T(t) = 0.2 - 0.8 (t + 0.5). At x + p = 1.0 (p = 1.5):
%   T = -1.0, phi = 0.8, gap = 1.8 = 0.8 p^2 (= phi''/2 p^2 exactly, a parabola).
%
% GREYSCALE: chord, tangent and graph differ by dash pattern and direct labels.
\definecolor{okabeblue}{HTML}{0072B2}

\begin{tikzpicture}[
    x=1.0cm, y=0.9cm,
    >={Latex[length=2.0mm]},
    font=\small,
    lab/.style={font=\scriptsize, inner sep=1.5pt},
    panelcap/.style={font=\small, align=center, text width=56mm, anchor=north},
  ]
% ------------------------------------------------------------- left: chord
\begin{scope}
  \draw[->, black!55] (-2.1,0) -- (2.2,0);
  \draw[->, black!55] (0,-0.25) -- (0,2.25);
  \node[lab, right] at (2.2,0) {$\mathbf{x}$};
  \draw[very thick, domain=-1.7:1.7, smooth, samples=50, variable=\t]
    plot ({\t},{0.8*\t*\t});
  \draw[okabeblue, thick] (-1.5,1.8) -- (1.0,0.8);
  \filldraw[okabeblue] (-1.5,1.8) circle (1.6pt);
  \filldraw[okabeblue] (1.0,0.8) circle (1.6pt);
  \draw[densely dotted] (-1.5,0) -- (-1.5,1.8);
  \draw[densely dotted] (1.0,0) -- (1.0,0.8);
  \node[lab, below] at (-1.5,0) {$\mathbf{x}_a$};
  \node[lab, below] at (1.0,0) {$\mathbf{x}_b$};
  \node[lab, below] at (-0.25,0) {$\mathbf{x}_c$};
  \draw[dashed] (-0.25,0) -- (-0.25,1.3);
  \filldraw (-0.25,1.3) circle (1.6pt);
  \filldraw (-0.25,0.05) circle (1.6pt);
  \node[lab, anchor=west] at (-0.1,0.45) {$\phi(\mathbf{x}_c)$};
  \node[lab, okabeblue, anchor=south west] at (-2.0,2.4)
    {chord value $\theta\phi(\mathbf{x}_a)+(1-\theta)\phi(\mathbf{x}_b)$};
  \node[panelcap] at (0,-0.5) {\textbf{chord above the graph}\\[-1pt] inequality \textcircled{\footnotesize 1}};
\end{scope}
% ------------------------------------------------------------- right: tangent
\begin{scope}[xshift=5.8cm]
  \draw[->, black!55] (-2.1,0) -- (2.2,0);
  \draw[->, black!55] (0,-0.4) -- (0,2.4);
  \node[lab, right] at (2.2,0) {$\mathbf{x}$};
  \draw[very thick, domain=-1.7:1.7, smooth, samples=50, variable=\t]
    plot ({\t},{0.8*\t*\t});
  % tangent at x = -0.5: T(t) = 0.2 - 0.8 (t + 0.5)
  \draw[okabeblue, thick, dashed, domain=-1.7:1.75, variable=\t]
    plot ({\t},{0.2-0.8*(\t+0.5)});
  \filldraw (-0.5,0.2) circle (1.6pt);
  % 2026-10-06: the label sat between the parabola, the tangent and the axis;
  % dropped to the axis like x_a, x_b, x_c in the chord panel.
  \draw[densely dotted] (-0.5,0.2) -- (-0.5,0);
  \node[lab, below] at (-0.5,0) {$\mathbf{x}$};
  \draw[densely dotted] (1.0,0.8) -- (1.0,-1.0);
  \filldraw (1.0,0.8) circle (1.6pt);
  \filldraw[okabeblue] (1.0,-1.0) circle (1.6pt);
  \draw[<->] (1.18,-1.0) -- (1.18,0.8);
  \node[lab, anchor=west] at (1.22,-0.55) {gap $\ge 0$};
    \node[lab, okabeblue, anchor=north east] at (1.0,-1.1) {$\phi(\mathbf{x}) + \nabla\phi(\mathbf{x})^{\top}\mathbf{p}$};
  \node[lab, anchor=south] at (1.0,0.95) {$\phi(\mathbf{x}+\mathbf{p})$};
  \node[panelcap] at (0,-1.75) {\textbf{tangent line below the graph}\\[-1pt] inequality \textcircled{\footnotesize 2}};
\end{scope}
\end{tikzpicture}
